\section*{Capítulo \ref{ch:b1:rate-laws} --- Cinética química: leyes de velocidad}

\begin{solution}{exo:b1:rate-laws:1}
$v = \frac14 \times 2.0 \times 10^{-4} = \qty{5.0e-5}{mol.L^{-1}.s^{-1}}$;
el dioxígeno se forma a la misma velocidad, \qty{5.0e-5}{}; \ce{N2O5}
desaparece a
$2v = \qty{1.0e-4}{mol.L^{-1}.s^{-1}}$.
\end{solution}

\begin{solution}{exo:b1:rate-laws:2}
Orden 0: \unit{mol.L^{-1}.s^{-1}}; orden 1: \unit{s^{-1}}; orden 2:
\unit{L.mol^{-1}.s^{-1}}; orden 3: \unit{L^2.mol^{-2}.s^{-1}}.
\end{solution}

\begin{solution}{exo:b1:rate-laws:3}
$t_{1/2} = \ln2/k = 0.693/(3.0 \times 10^{-3}) = \qty{231}{s}$. Al cabo de
\qty{600}{s}, $\eu^{-kt} = \eu^{-1.8} = 0.165$: queda un 16.5\,\%.
\end{solution}

\begin{solution}{exo:b1:rate-laws:4}
Un \omterm{def:b1:rate-laws:half-life}{tiempo de semirreacción} proporcional a $1/a_0$ indica orden 2. Si al
dividir $a_0$ por dos el \omterm{def:b1:rate-laws:half-life}{tiempo de semirreacción} se divide por dos,
$t_{1/2} \propto a_0$: orden 0.
\end{solution}

\begin{solution}{exo:b1:rate-laws:5}
$[\ce{A}]$ no es lineal en $t$ (pérdidas: 0.0259, 0.0192, 0.0142, 0.0106).
$\ln[\ce{A}] = -2.303$, $-2.602$, $-2.902$, $-3.202$, $-3.503$: disminuye
0.300 cada \qty{10}{min}, una recta. Orden 1, $k =
\qty{0.0300}{min^{-1}}$.
\end{solution}

\begin{solution}{exo:b1:rate-laws:6}
Duplicar $[\ce{A}]_0$ multiplica $v_0$ por $4 = 2^2$: orden 2 respecto de
\ce{A}. Duplicar $[\ce{B}]_0$ la multiplica por 2: orden 1 respecto de
\ce{B}. $v =
k[\ce{A}]^2[\ce{B}]$ con $k = 2.0 \times 10^{-6}/(0.010^2 \times 0.010) =
\qty{2.0}{L^2.mol^{-2}.s^{-1}}$.
\end{solution}

\begin{solution}{exo:b1:rate-laws:7}
\ce{B} está en un exceso de cien veces: su concentración se mantiene en
\qty{0.50}{mol/L}, y $v = k'[\ce{A}]$ con $k' = k[\ce{B}]_0$. Por tanto,
$k = 0.012/0.50 = \qty{0.024}{L.mol^{-1}.s^{-1}}$.
\end{solution}

\begin{solution}{exo:b1:rate-laws:8}
$E_a = R\ln5/(1/300 - 1/320) = 8.314 \times 1.609/(2.083 \times 10^{-4}) =
\qty{64.2}{kJ/mol}$. A \qty{310}{K}: $\ln k = \ln(1.0 \times 10^{-3}) +
\frac{E_a}{R}\big(\frac{1}{300} - \frac{1}{310}\big) = -6.908 + 0.831$,
$k = \qty{2.3e-3}{s^{-1}}$.
\end{solution}

\begin{solution}{exo:b1:rate-laws:9}
Aquí $a = 2$: $1/[\ce{A}] = 1/a_0 + 2kt$, $t_{1/2} = 1/(2k a_0) = 1/(2
\times 0.50 \times 0.020) = \qty{50}{s}$. A \qty{100}{s}: $1/[\ce{A}] = 50
+ 100 = \qty{150}{L/mol}$, $[\ce{A}] = \qty{6.7e-3}{mol/L}$.
\end{solution}

\begin{solution}{exo:b1:rate-laws:10}
$[\ce{A}] = a_0(1 - f) = a_0\eu^{-akt_f}$ da $t_f = -\ln(1-f)/(ak)$ y, en
\omterm{def:b1:rate-laws:half-life}{tiempos de semirreacción}, $t_f/t_{1/2} = -\ln(1-f)/\ln2$: 1 para el
50\,\%, 2 para el 75\,\%, 3.32 para el 90\,\% y 9.97 para el 99.9\,\%:
unos diez \omterm{def:b1:rate-laws:half-life}{tiempos de semirreacción} para que la reacción haya
``terminado''.
\end{solution}

\begin{solution}{exo:b1:rate-laws:11}
$m(t) = 10 - 0.40\,t$ (\unit{mg}, \unit{h}): orden 0. \omterm{def:b1:rate-laws:half-life}{Tiempo de
semirreacción}: $10/(2 \times 0.40) = \qty{12.5}{h}$; vacío al cabo de
\qty{25}{h}. La dosis liberada por hora no depende de lo que queda: un
aporte constante, que es la finalidad de un parche.
\end{solution}

\begin{solution}{exo:b1:rate-laws:12}
$E_a = R\ln2/(1/298.15 - 1/308.15) = 8.314 \times 0.693/(1.088 \times
10^{-4}) = \qty{52.9}{kJ/mol}$. Entre 273.15 y \qty{283.15}{K}:
$\exp\big(\frac{52900}{8.314}(\frac{1}{273.15} - \frac{1}{283.15})\big) =
2.3$: la misma \omterm{def:b1:rate-laws:activation-energy}{energía de activación} da un factor mayor a temperatura más
baja.
\end{solution}

\begin{solution}{pb:b1:rate-laws:1}
\textbf{1.} Pérdidas de 8.6, 7.9, 7.2 y 6.5 puntos: no son constantes,
luego no es de orden 0.
\textbf{2.} $\ln$: 4.605, 4.515, 4.425, 4.335, 4.245; disminuye 0.090
cada 10 días: una recta, orden 1.
\textbf{3.} $k_{60} = 0.090/10 = \qty{9.0e-3}{d^{-1}}$.
\textbf{4.} $t_{1/2} = \ln2/k_{60} = \qty{77}{d}$.
\textbf{5.} $100\,\eu^{-0.54} = 58.3\,\%$.
\textbf{6.} $t_{90} = \ln(10/9)/k_{60} = 0.1054/0.0090 = \qty{11.7}{d}$.
\textbf{7.} Queda un 90\,\%: $\eu^{-kt_{90}} = 0.9$, de modo que $t_{90} = \ln(10/9)/k$.
\textbf{8.} $t_{90}/t_{1/2} = \ln(10/9)/\ln2 = 0.152$.
\textbf{9.} En el orden 1, la fracción restante solo depende de $kt$, no
de la cantidad inicial.
\textbf{10.} A \qty{25}{\celsius}, la degradación tarda años; calentar la
acelera (Arrhenius) para poder medirla en semanas.
\textbf{11.} En el orden 2, el \omterm{def:b1:rate-laws:half-life}{tiempo de semirreacción} es $1/(ak\,a_0)$:
un comprimido del doble de dosis se degradaría el doble de rápido en
términos relativos y se conservaría menos tiempo.
\textbf{12.} $1/T$ (\unit{K^{-1}}): \num{3.1934e-3}, \num{3.0945e-3},
\num{3.0017e-3}; $\ln k$: $-6.669$, $-5.661$, $-4.711$.
\textbf{13.} Pendiente $= (-4.711 + 6.669)/(3.0017 \times 10^{-3} - 3.1934 \times
10^{-3}) = \qty{-1.021e4}{K}$.
\textbf{14.} $E_a = 8.314 \times 1.021 \times 10^4 = \qty{85}{kJ/mol}$.
\textbf{15.} Valor predicho: $\ln k_{50} = -4.711 - 1.021 \times 10^4 \times (3.0945 -
3.0017) \times 10^{-3} = -5.658$; medido, $-5.661$: sobre la recta.
\textbf{16.} $A = k_{60}\,\eu^{E_a/RT} = 0.0090\,\eu^{10215/333.15} =
\qty{1.9e11}{d^{-1}}$.
\textbf{17.} $\exp\big(1.021 \times 10^4(\frac{1}{298.15} -
\frac{1}{308.15})\big) = 3.0$: más que el factor 2 de la regla, porque
$E_a$ es mayor que \qty{53}{kJ/mol}.
\textbf{18.} $k_{30} = k_{60}\exp\big(-1.021 \times 10^4(\frac{1}{303.15} -
\frac{1}{333.15})\big) = \qty{4.3e-4}{d^{-1}}$.
\textbf{19.} $t_{90} = 0.1054/4.33 \times 10^{-4} = \qty{243}{d}$, 8.0
meses.
\textbf{20.} $1 - \eu^{-3 \times 0.0090} = 2.7\,\%$.
\textbf{21.} La vida útil disminuye rápidamente con la temperatura (8
meses a \qty{30}{\celsius} frente a 14 a \qty{25}{\celsius}): la fecha de
caducidad solo es válida por debajo de \qty{25}{\celsius}.
\textbf{22.} $k_{25} = k_{60}\exp\big(-1.021 \times 10^4(\frac{1}{298.15} -
\frac{1}{333.15})\big) = \qty{2.46e-4}{d^{-1}}$.
\textbf{23.} $t_{90} = 0.1054/2.46 \times 10^{-4} = \qty{428}{d}$, es
decir, \textbf{14 meses}: la fecha impresa en la caja, obtenida a partir
de seis semanas de medidas.
\end{solution}
